\begin{proof}[Proof of \cref{obj:prop:bounded-outcome-gaussian-degeneracy}]
\begin{enumerate}
\item Fix a model \(m\) in the displayed Gaussian class, and write \(P\) for its observed-data law. Define
\[
\eta=T-g_0(X),
\qquad
\xi=Y-q_0(X)-\theta_0(P)\eta .
\]
The conditional mean restrictions in \cref{obj:def:gaussian-class,obj:def:plm-model} give
\[
E_P[\xi\mid X,T]=0,
\]
and \(q_0(X)\) and \(\theta_0(P)\eta\) are measurable with respect to \((X,T)\). Hence conditional-expectation linearity yields
\[
E_P[Y\mid X,T]
=
q_0(X)+\theta_0(P)\eta
\qquad P\text{-a.s.}
\]
The bounded-outcome condition gives \(-C_q\le Y\le C_q\) \(P\)-a.s. Conditioning the two inequalities on \((X,T)\) gives
\[
-C_q\le E_P[Y\mid X,T]\le C_q
\qquad P\text{-a.s.},
\]
and therefore
\[
\bigl|q_0(X)+\theta_0(P)\eta\bigr|\le C_q
\qquad P\text{-a.s.}
\]
Conditioning the same bounded-outcome inequalities on \(X\), and using
\[
E_P[Y\mid X]=q_0(X),
\]
gives
\[
|q_0(X)|\le C_q
\qquad P\text{-a.s.}
\]
Combining the two displayed bounds,
\[
|\theta_0(P)\eta|
=
\bigl|(q_0(X)+\theta_0(P)\eta)-q_0(X)\bigr|
\le
\bigl|q_0(X)+\theta_0(P)\eta\bigr|+|q_0(X)|
\le 2C_q
\qquad P\text{-a.s.}
\]
% lean: CausalSmith.Stat.SaPlmCumulantConverse.gaussianClass_theta_eq_zero

\item Suppose, for contradiction, that \(\theta_0(P)\ne0\), and define
\[
\Lambda_\theta=\frac{2C_q}{|\theta_0(P)|}\in\mathbb R .
\]
The previous step implies
\[
|\eta|\le \Lambda_\theta
\qquad P\text{-a.s.},
\]
so
\[
P\{\eta>\Lambda_\theta\}=0.
\]
The Gaussian membership in \cref{obj:def:gaussian-class} gives \(\eta\sim N(0,\sigma^2)\), and the parameter record in \cref{obj:def:model-parameters} gives \(\sigma>0\). Thus the variance is nonzero and
\[
N(0,\sigma^2)\bigl((\Lambda_\theta,\infty)\bigr)=0.
\]
For a Gaussian law with positive variance, Lebesgue measure is absolutely continuous with respect to that Gaussian law. Hence the last display implies
\[
\lambda\bigl((\Lambda_\theta,\infty)\bigr)=0.
\]
But
\[
\lambda\bigl((\Lambda_\theta,\infty)\bigr)=\infty .
\]
This contradiction proves
\[
\theta_0(P)=0 .
\]
% lean: CausalSmith.Stat.SaPlmCumulantConverse.gaussianClass_theta_eq_zero

\item Now fix deterministic code sequences \(\widehat g,\widehat q\) and assume the fixed-code Gaussian class in the statement is inhabited. Consider the estimator
\[
\widetilde\theta_0(O_1,\ldots,O_n;\widehat g_n,\widehat q_n)=0 .
\]
For every model in that fixed-code class, the first part gives \(\theta_0(P)=0\), and hence
\[
E_{P^{\otimes n}}\!\left[
\{\widetilde\theta_0(O_1,\ldots,O_n;\widehat g_n,\widehat q_n)-\theta_0(P)\}^2
\right]=0 .
\]
Taking the supremum over the fixed-code class and then the infimum over all estimators in \cref{obj:def:minimax-risks} gives
\[
\mathfrak R_n^{\mathrm G}(p;\widehat g,\widehat q)\le0.
\]
Since this risk is nonnegative by definition,
\[
\mathfrak R_n^{\mathrm G}(p;\widehat g,\widehat q)=0 .
\]
% lean: CausalSmith.Stat.SaPlmCumulantConverse.bounded_outcome_gaussian_degeneracy

\item It remains to evaluate the generalized-quantile criterion for the same zero estimator. Since \(\gamma\in(1/2,1)\),
\[
0<1-\gamma<1.
\]
For every model in the fixed-code Gaussian class,
\[
\bigl|\widetilde\theta_0(O_1,\ldots,O_n;\widehat g_n,\widehat q_n)-\theta_0(P)\bigr|=0
\qquad P^{\otimes n}\text{-a.s.}
\]
Thus the law of this loss is the point mass at zero. Its distribution function is \(0\) for negative thresholds and \(1\) at every nonnegative threshold, so the generalized lower \((1-\gamma)\)-quantile of \cref{obj:def:generalized-quantile} is
\[
Q_{P,1-\gamma}(0)=0 .
\]
Consequently the supremum of the fixed-code quantile loss for the zero estimator is \(0\), and the minimax infimum is at most \(0\). Nonnegativity in \cref{obj:def:minimax-risks} gives
\[
\mathfrak M_{n,1-\gamma}^{\mathrm G}(p;\widehat g,\widehat q)=0 .
\]
% lean: CausalSmith.Stat.SaPlmCumulantConverse.generalizedQuantile_zero
\end{enumerate}
\end{proof}
